\section{Simulation studies}
\label{simuStudy}

We conducted simulation analyses to assess the performance of the proposed
fixed-threshold estimation scheme with two biomarkers. In contrast to a model
that estimates a separate cut-point and constrains the biomarker coefficients
through angular coordinates, the present formulation fixes the threshold at
zero and the intercept of the single-index score at one. Thus, for subject
$i$, the biomarker-defined subgroup is
\[
    G_i(\boldsymbol\gamma)
    = I\{1+\mathbf{z}_i^{T}\boldsymbol\gamma\geq 0\},
\]
and event times were generated from the exponential model with hazard
\begin{equation}
\label{eq:simModelFixedThreshold}
    h_i(t)
    = \exp\left\{
        \beta_1 a_i
        + \beta_2 G_i(\boldsymbol\gamma)
        + \beta_3 a_iG_i(\boldsymbol\gamma)
      \right\},
\end{equation}
where $a_i$ is the treatment indicator. The coefficient $\beta_1$ is the
treatment effect in the index-negative subgroup, $\beta_2$ is the main effect
of the index-positive subgroup, and $\beta_3$ is the treatment-by-subgroup
interaction. Consequently, the treatment hazard ratios in the index-negative
and index-positive subgroups are $\exp(\beta_1)$ and
$\exp(\beta_1+\beta_3)$, respectively.

We first examined a grid of regression parameters with sample sizes
$n=600$, $1000$, and $1400$. Patients were independently assigned to the
treatment or control arm with probability 0.5. The two biomarkers were
generated independently from standard normal distributions, and their true
index coefficients were fixed at
$\boldsymbol\gamma=(1.5,1.8)^{T}$. We considered
$\exp(\beta_1)\in\{1,2\}$,
$\exp(\beta_2)\in\{2.5,2.8\}$, and
$\exp(\beta_3)\in\{2.5,2.8\}$. Including $\beta_1=0$ allowed us to examine
whether the biomarker-defined subgroup could be estimated in the absence of
an overall treatment effect. The baseline hazard was one and independent
administrative censoring times were generated from $U(0.4,3)$. Across the
parameter grid, the empirical censoring proportions ranged from approximately
6.9\% to 10.1\%.

For each parameter combination, we obtained 500 successful simulated data
sets, organized as 50 computational replicates of 10 data sets. A generated
data set was eligible when each of the four treatment-by-subgroup cells
contained at least 60 patients. We used the unpenalized model
($\lambda=0$), initialized $\boldsymbol\gamma$ at $(1,1.3)^{T}$, and assigned
independent $N(0.5,1)$ priors to its two components. The normal random-walk
proposal standard deviation for each component of $\boldsymbol\gamma$ was
0.5. Each MCMC run used 10,000 burn-in iterations and retained 20,000
posterior draws after thinning by 10, for a total of 210,000 iterations.
Posterior means were used as point estimates, and empirical coverage was
calculated from equal-tailed 95\% credible intervals.

Table~\ref{tab:betaGridSummary} summarizes the range of mean absolute bias,
root mean squared error (RMSE), and empirical coverage across the eight
regression-parameter combinations at each sample size. The signed biases were
small throughout: the largest absolute mean bias among all coefficients and
scenarios was 0.014. The mean absolute biases and RMSEs generally decreased as
sample size increased. This pattern was especially clear for the index
coefficients. For example, the mean absolute bias of $\hat\gamma_1$ decreased
from 0.045--0.053 at $n=600$ to 0.017--0.021 at $n=1400$, while the
corresponding range for $\hat\gamma_2$ decreased from 0.049--0.057 to
0.020--0.023. Coverage probabilities remained close to the nominal 95\%
level, ranging from 0.930 to 0.966 over the full grid. All 12,000 fits in this
grid converged.

\begin{table}[!htbp]
\caption[Simulation performance over the regression-parameter grid]{Ranges of
mean absolute bias, RMSE, and empirical coverage over the eight regression-
parameter combinations at each sample size. Each parameter combination was
represented by 500 successfully fitted simulated data sets.}
\label{tab:betaGridSummary}
\centering
\scriptsize
\begin{tabular}{clccc}
\toprule
$n$ & Parameter & Mean absolute bias & RMSE & Coverage probability \\
\midrule
600  & $\beta_1$  & 0.123--0.124 & 0.154--0.156 & 0.958--0.962 \\
600  & $\beta_2$  & 0.114--0.115 & 0.143--0.145 & 0.940--0.950 \\
600  & $\beta_3$  & 0.147--0.154 & 0.185--0.194 & 0.960--0.966 \\
600  & $\gamma_1$ & 0.045--0.053 & 0.062--0.072 & 0.932--0.942 \\
600  & $\gamma_2$ & 0.049--0.057 & 0.072--0.086 & 0.930--0.948 \\
\addlinespace
1000 & $\beta_1$  & 0.095--0.103 & 0.123--0.131 & 0.934--0.946 \\
1000 & $\beta_2$  & 0.086--0.087 & 0.106--0.108 & 0.954--0.960 \\
1000 & $\beta_3$  & 0.123--0.125 & 0.156--0.159 & 0.934--0.944 \\
1000 & $\gamma_1$ & 0.025--0.030 & 0.033--0.040 & 0.946--0.954 \\
1000 & $\gamma_2$ & 0.026--0.031 & 0.035--0.043 & 0.956--0.964 \\
\addlinespace
1400 & $\beta_1$  & 0.083--0.085 & 0.104--0.107 & 0.952--0.956 \\
1400 & $\beta_2$  & 0.071--0.071 & 0.089--0.090 & 0.942--0.954 \\
1400 & $\beta_3$  & 0.102--0.103 & 0.127--0.128 & 0.948--0.960 \\
1400 & $\gamma_1$ & 0.017--0.021 & 0.024--0.029 & 0.940--0.946 \\
1400 & $\gamma_2$ & 0.020--0.023 & 0.028--0.034 & 0.938--0.952 \\
\bottomrule
\end{tabular}
\end{table}

We additionally examined a fixed parameter setting over a finer sequence of
sample sizes. In this study,
$\boldsymbol\beta=[\log(1.5),\log(1.2),\log(2)]^{T}$ and
$\boldsymbol\gamma=(1.5,1.8)^{T}$, with 500 successfully fitted data sets for
each of $n=600$, $800$, $1000$, $1200$, and $1400$. The results are displayed
in Table~\ref{tab:fixedParameterSummary}. The signed biases approached zero
with increasing sample size, and both the mean absolute biases and RMSEs
decreased for every parameter. For example, the mean absolute biases of
$\hat\gamma_1$ and $\hat\gamma_2$ decreased from 0.176 and 0.194 at $n=600$
to 0.093 and 0.095 at $n=1400$. The corresponding RMSEs decreased from 0.237
and 0.258 to 0.127 and 0.133. Coverage probabilities ranged from 0.934 to
0.970 and showed no systematic deterioration as the credible intervals
narrowed. The convergence rate was 95.4\% at $n=600$, where 524 fits were
attempted to obtain 500 successful fits; it increased to 99.8\% at $n=800$
and was 100\% for $n\geq1000$.

\begin{table}[!htbp]
\caption[Fixed-parameter simulation performance by sample size]{Simulation
performance for
$\boldsymbol\beta=[\log(1.5),\log(1.2),\log(2)]^{T}$ and
$\boldsymbol\gamma=(1.5,1.8)^{T}$. Each sample size is represented by 500
successfully fitted simulated data sets.}
\label{tab:fixedParameterSummary}
\centering
\scriptsize
\begin{tabular}{clrrrr}
\toprule
$n$ & Parameter & Bias & Mean absolute bias & RMSE & Coverage \\
\midrule
600 & $\beta_1$  & -0.015 & 0.125 & 0.158 & 0.962 \\
600 & $\beta_2$  &  0.006 & 0.116 & 0.144 & 0.970 \\
600 & $\beta_3$  &  0.022 & 0.153 & 0.192 & 0.956 \\
600 & $\gamma_1$ & -0.043 & 0.176 & 0.237 & 0.946 \\
600 & $\gamma_2$ & -0.067 & 0.194 & 0.258 & 0.960 \\
\addlinespace
800 & $\beta_1$  & -0.014 & 0.117 & 0.146 & 0.938 \\
800 & $\beta_2$  & -0.006 & 0.104 & 0.131 & 0.934 \\
800 & $\beta_3$  &  0.019 & 0.139 & 0.172 & 0.962 \\
800 & $\gamma_1$ & -0.015 & 0.150 & 0.209 & 0.934 \\
800 & $\gamma_2$ & -0.031 & 0.163 & 0.221 & 0.942 \\
\addlinespace
1000 & $\beta_1$  & -0.003 & 0.099 & 0.126 & 0.940 \\
1000 & $\beta_2$  &  0.002 & 0.089 & 0.111 & 0.956 \\
1000 & $\beta_3$  &  0.010 & 0.125 & 0.156 & 0.938 \\
1000 & $\gamma_1$ & -0.014 & 0.119 & 0.160 & 0.952 \\
1000 & $\gamma_2$ & -0.013 & 0.122 & 0.166 & 0.960 \\
\addlinespace
1200 & $\beta_1$  &  0.005 & 0.089 & 0.111 & 0.964 \\
1200 & $\beta_2$  &  0.001 & 0.078 & 0.099 & 0.958 \\
1200 & $\beta_3$  & -0.001 & 0.106 & 0.133 & 0.960 \\
1200 & $\gamma_1$ & -0.004 & 0.102 & 0.141 & 0.954 \\
1200 & $\gamma_2$ & -0.009 & 0.110 & 0.155 & 0.958 \\
\addlinespace
1400 & $\beta_1$  & -0.000 & 0.083 & 0.104 & 0.944 \\
1400 & $\beta_2$  & -0.000 & 0.071 & 0.088 & 0.952 \\
1400 & $\beta_3$  &  0.002 & 0.099 & 0.124 & 0.960 \\
1400 & $\gamma_1$ &  0.003 & 0.093 & 0.127 & 0.942 \\
1400 & $\gamma_2$ & -0.006 & 0.095 & 0.133 & 0.948 \\
\bottomrule
\end{tabular}
\end{table}

% FIGURE NEEDED HERE: A five-panel or faceted line plot of mean absolute bias
% and RMSE against n for beta_1, beta_2, beta_3, gamma_1, and gamma_2 in the
% fixed-parameter study. The plot should be made from the values in
% Table~\ref{tab:fixedParameterSummary} and should emphasize the decrease in
% estimation error with increasing sample size.
\begin{figure}[!htbp]
\centering
\fbox{\parbox{0.92\textwidth}{\centering\textbf{Figure needed:} mean absolute
bias and RMSE versus sample size for the five parameters in the
fixed-parameter simulation.}}
\caption[Estimation error by sample size]{Mean absolute bias and RMSE by
sample size in the fixed-parameter simulation. Separate curves are needed for
$\beta_1$, $\beta_2$, $\beta_3$, $\gamma_1$, and $\gamma_2$.}
\label{fig:simBiasByN}
\end{figure}

Taken together, the simulations indicate that fixing the index intercept and
threshold yields stable estimation of both the regression and biomarker-index
coefficients in the settings examined. The strongest evidence is the
consistent reduction in absolute bias and RMSE as $n$ increases, together
with coverage probabilities that remain near 0.95. The results with
$\beta_1=0$ also show that the model can recover the subgroup and interaction
parameters when there is no overall treatment effect in the index-negative
subgroup.

\section{Data Analysis}
\label{analysis}

For application, we used data from a phase III randomized trial of patients
with locally advanced or metastatic pancreatic cancer
\cite{doi:10.1200/JCO.2006.07.9525}. Patients were randomly assigned to
receive gemcitabine plus erlotinib or gemcitabine alone. Protein expression
levels measured from pretreatment plasma samples have previously been used to
study prognostic and treatment-predictive biomarkers in this trial
\cite{ShultzDavidB2016ANBP,HeYe2018AstC}. In particular, HER2 has been
reported as a prognostic biomarker, whereas the combination of CA19-9 and Axl
has been investigated as a potential definition of erlotinib sensitivity
subsets.

The analysis data contained 569 patients, including 285 in the erlotinib arm
and 284 in the control arm; 485 events were observed. For each protein, values
that could not be converted to numeric form were treated as missing and
replaced by the observed column mean. The biomarker values were then
log-transformed and standardized. We focused on HER2, CA19-9, and Axl. We
first fitted the model with all three biomarkers and then fitted each of the
three two-biomarker combinations for comparison.

For every model we used four MCMC chains. Each chain discarded 50,000 burn-in
iterations and retained 50,000 posterior samples after thinning by 10, for
550,000 iterations per chain. Independent $N(0,1)$ priors were assigned to
the components of $\boldsymbol\gamma$, and no lasso penalty was applied
($\lambda=0$). For the three-biomarker model, the starting value was
$(-0.6,0.5,0.5)^{T}$ for HER2, CA19-9, and Axl, respectively. Posterior means
and equal-tailed 95\% credible intervals are reported below. The split
$\widehat R$ values for the three-biomarker model ranged from 1.001 to 1.006,
and the effective sample sizes ranged from 589 to 1801. Across the pairwise
models, the largest split $\widehat R$ was 1.014. These diagnostics and the
trace plots did not indicate gross between-chain disagreement, although the
moderate effective sample sizes for several $\gamma$ coefficients reflect
substantial autocorrelation.

% FIGURE NEEDED HERE: Four-chain trace plots and autocorrelation plots for all
% beta and gamma parameters in the HER2/CA19-9/Axl analysis. Existing source
% files: results/biomarker-mcmc-trial3/convergence-traces.pdf and
% results/biomarker-mcmc-trial3/convergence-autocorrelation.pdf.
\begin{figure}[!htbp]
\centering
\fbox{\parbox{0.92\textwidth}{\centering\textbf{Figure needed:} four-chain
trace and autocorrelation plots for the HER2, CA19-9, and Axl model. Existing
plots are stored in \texttt{results/biomarker-mcmc-trial3/}.}}
\caption[MCMC diagnostics for the three-biomarker model]{Trace and
autocorrelation plots for the four MCMC chains in the HER2, CA19-9, and Axl
analysis.}
\label{fig:analysisConvergence}
\end{figure}

Table~\ref{tab:analysisEst} presents the posterior estimates for the
three-biomarker analysis. The estimated score was
\begin{equation}
\label{eq:threeBiomarkerScore}
    S_i
    = 1 - 0.176\,\mathrm{HER2}_i
        + 0.389\,\mathrm{CA19\mbox{-}9}_i
        - 0.648\,\mathrm{Axl}_i,
\end{equation}
with patients classified as index-positive when $S_i\geq0$. All three
biomarker credible intervals contained zero. The estimated
treatment-by-subgroup interaction was $\hat\beta_3=0.548$, with a 95\%
credible interval from $-1.093$ to $1.864$, and therefore did not provide
strong posterior evidence of treatment-effect heterogeneity.

\begin{table}[!htbp]
\caption[Posterior estimates for the three-biomarker analysis]{Posterior means,
posterior standard deviations, equal-tailed 95\% credible intervals, split
$\widehat R$, and effective sample sizes for the model containing HER2,
CA19-9, and Axl.}
\label{tab:analysisEst}
\centering
\scriptsize
\begin{tabular}{lrrrrrr}
\toprule
Parameter & Estimate & Posterior SD & Lower 95\% & Upper 95\% & Split $\widehat R$ & ESS \\
\midrule
$\beta_1$ (treatment)             & -0.670 & 0.726 & -1.998 & 0.778 & 1.005 & 1195 \\
$\beta_2$ (index-positive group)  &  0.026 & 0.543 & -0.529 & 0.702 & 1.003 & 1801 \\
$\beta_3$ (interaction)           &  0.548 & 0.768 & -1.093 & 1.864 & 1.006 & 1028 \\
$\gamma_{\mathrm{HER2}}$          & -0.176 & 0.565 & -1.058 & 1.191 & 1.006 & 686 \\
$\gamma_{\mathrm{CA19\mbox{-}9}}$&  0.389 & 0.565 & -0.897 & 1.306 & 1.001 & 942 \\
$\gamma_{\mathrm{Axl}}$           & -0.648 & 0.684 & -1.936 & 0.936 & 1.005 & 589 \\
\bottomrule
\end{tabular}
\end{table}

Using the posterior-mean score in Equation~\ref{eq:threeBiomarkerScore}, 41
patients were classified as index-negative and 528 as index-positive. In the
index-negative subgroup, median survival was 8.31 months among treated
patients and 6.46 months among controls; the treatment log-rank p-value was
0.114. In the index-positive subgroup, the corresponding median survival
times were 6.11 and 5.91 months, with p=0.089. The four-group log-rank test
gave p=0.047. When the treatment hazard ratio was calculated from every
posterior draw, its posterior mean and 95\% credible interval were 0.68
(0.14, 2.18) in the index-negative subgroup and 0.89 (0.68, 1.11) in the
index-positive subgroup. A conventional Cox model applied after fixing the
estimated subgroup gave an interaction hazard ratio of 1.47 (95\% confidence
interval 0.72--3.03; Wald p=0.293). Thus, although the descriptive survival
curves differed across the four groups, the three-biomarker analysis did not
show a precise treatment-by-subgroup interaction.

% FIGURE NEEDED HERE: Kaplan--Meier curves for the four treatment-by-subgroup
% groups under the HER2/CA19-9/Axl posterior-mean score. Existing source:
% results/km-from-mcmc/HER2_CA19-9_Axl/km-four-groups.pdf. A two-panel version
% may additionally use km-treatment-below-threshold.pdf and
% km-treatment-above-threshold.pdf from the same directory.
\begin{figure}[!htbp]
\centering
\fbox{\parbox{0.92\textwidth}{\centering\textbf{Figure needed:}
Kaplan--Meier curves for the four treatment-by-subgroup groups from the HER2,
CA19-9, and Axl model, with numbers at risk.}}
\caption[Kaplan--Meier curves for the three-biomarker model]{Kaplan--Meier
curves by treatment and posterior-mean biomarker subgroup for the HER2,
CA19-9, and Axl model.}
\label{fig:threeBiomarkerKM}
\end{figure}

For comparison, we also fitted the model using each pair of biomarkers.
Table~\ref{tab:pairwiseAnalysis} summarizes the regression coefficients. The
credible intervals for all three interactions contained zero. The pairwise
scores and their descriptive survival results are summarized in
Table~\ref{tab:pairwiseSurvival}. Of the three fits, the CA19-9/Axl model
produced the most distinct and usable partition. Its posterior-mean score was
\begin{equation}
\label{eq:caAxlScore}
    S_i=1+0.641\,\mathrm{CA19\mbox{-}9}_i-0.921\,\mathrm{Axl}_i.
\end{equation}
The 95\% credible intervals for both index coefficients contained zero:
$(-0.721,1.340)$ for CA19-9 and $(-2.080,0.347)$ for Axl. Seventy-four
patients were below the fixed threshold. Within this subgroup,
median survival was 8.31 months in the treatment arm and 5.65 months in the
control arm, and the treatment log-rank p-value was 0.010. Among the 495
index-positive patients, median survival was 6.08 and 5.95 months,
respectively, with p=0.284. The draw-by-draw posterior treatment hazard ratio
was 0.77 (95\% credible interval 0.015--0.996) below the threshold and 0.91
(0.73--1.11) above it. The post-classification Cox interaction hazard ratio
was 1.77 (95\% confidence interval 1.01--3.08; Wald p=0.045).

\begin{table}[!htbp]
\caption[Posterior regression estimates for the pairwise analyses]{Posterior
means and equal-tailed 95\% credible intervals for the regression coefficients
in the three pairwise biomarker models.}
\label{tab:pairwiseAnalysis}
\centering
\scriptsize
\begin{tabular}{lccc}
\toprule
Biomarkers & $\beta_1$ & $\beta_2$ & $\beta_3$ \\
\midrule
HER2, CA19-9
& 0.076 $(-1.344,1.487)$
& 0.238 $(-0.737,1.041)$
& -0.305 $(-1.743,1.224)$ \\
HER2, Axl
& -0.328 $(-2.869,1.107)$
& 0.109 $(-2.110,0.963)$
& 0.120 $(-1.418,2.691)$ \\
CA19-9, Axl
& -0.859 $(-4.211,-0.004)$
& -0.115 $(-4.322,0.573)$
& 0.758 $(-0.172,4.022)$ \\
\bottomrule
\end{tabular}
\end{table}

\begin{table}[!htbp]
\caption[Subgroup sizes and exploratory survival comparisons]{Subgroup sizes,
median survival times, and exploratory treatment comparisons based on the
posterior-mean pairwise scores. The Cox interaction tests condition on
subgroups estimated from the same data.}
\label{tab:pairwiseSurvival}
\centering
\scriptsize
\begin{tabular}{lrrrrrr}
\toprule
& \multicolumn{2}{c}{Subgroup size}
& \multicolumn{2}{c}{Treated/control median}
& \multicolumn{1}{c}{Treatment log-rank p}
& \multicolumn{1}{c}{Interaction} \\
\cmidrule(lr){2-3}\cmidrule(lr){4-5}
Biomarkers & Below & Above & Below & Above & Below / Above & HR (Wald p) \\
\midrule
HER2, CA19-9 & 18 & 551 & 4.24/7.36 & 6.37/5.91 & 0.122 / 0.012 & 0.30 (0.017) \\
HER2, Axl    &  0 & 569 & ---         & 6.37/5.91 & --- / 0.035   & not estimable \\
CA19-9, Axl  & 74 & 495 & 8.31/5.65 & 6.08/5.95 & 0.010 / 0.284 & 1.77 (0.045) \\
\bottomrule
\end{tabular}
\end{table}

The remaining pairwise fits require more caution. The HER2/CA19-9 score
was $1+0.480\,\mathrm{HER2}+0.011\,\mathrm{CA19\mbox{-}9}$; it classified
only 18 patients as index-negative, and the posterior mean of the CA19-9
coefficient was close to zero. The HER2/Axl score was
$1+0.076\,\mathrm{HER2}-0.149\,\mathrm{Axl}$. Its posterior-mean value
classified all 569 patients as index-positive, so an interaction was not
estimable from the resulting partition. These outcomes illustrate that
convergence of the MCMC chains does not by itself guarantee a clinically
useful or adequately sized subgroup.

% FIGURE NEEDED HERE: A two-panel Kaplan--Meier display for the CA19-9/Axl
% model, showing treatment within the below-threshold and above-threshold
% subgroups, with numbers at risk. Existing sources:
% results/km-from-mcmc/CA19-9_Axl/km-treatment-below-threshold.pdf and
% results/km-from-mcmc/CA19-9_Axl/km-treatment-above-threshold.pdf.
\begin{figure}[!htbp]
\centering
\fbox{\parbox{0.92\textwidth}{\centering\textbf{Figure needed:}
treatment-specific Kaplan--Meier curves within the below- and above-threshold
subgroups from the CA19-9/Axl model, with numbers at risk.}}
\caption[Kaplan--Meier curves for the CA19-9/Axl model]{Kaplan--Meier curves
by treatment within the index-negative and index-positive subgroups from the
CA19-9/Axl posterior-mean score.}
\label{fig:caAxlKM}
\end{figure}

All log-rank and conventional Cox p-values in this section are exploratory.
The biomarker coefficients and subgroup assignments were estimated from the
same observations used for the survival comparisons, so these tests do not
account for score-selection uncertainty and may be anti-conservative. The
posterior intervals, subgroup sizes, and stability of the estimated score
should therefore be considered together rather than interpreting the
post-classification p-values in isolation.

\section{Discussion}

We developed a Bayesian fixed-threshold single-index model for identifying
treatment-sensitive subgroups in survival data with multiple biomarkers. The
model uses the score $1+\mathbf{z}_i^{T}\boldsymbol\gamma$ and fixes its
threshold at zero. This parameterization removes the need to estimate a
separate cut-point or represent the biomarker coefficients through angular
coordinates. It also anchors the scale of the index by fixing its intercept,
thereby avoiding the cut-point--coefficient rescaling ambiguity that arises
when both quantities are freely estimated. A Metropolis-within-Gibbs
algorithm is used to sample the regression coefficients and the biomarker
index coefficients from the posterior distribution.

The simulation results suggest that the proposed formulation can recover the
regression and index coefficients under the conditions examined. Across 24
sample-size and regression-parameter scenarios, signed biases were small and
empirical coverage was generally close to 95\%. The additional fixed-
parameter study showed a consistent decrease in mean absolute bias and RMSE
from $n=600$ to $n=1400$. This improvement was especially pronounced for
$\boldsymbol\gamma$, which is important because uncertainty in the index
coefficients directly affects subgroup membership. The simulations with
$\beta_1=0$ further indicate that the method can estimate a biomarker-defined
interaction even when there is no treatment effect in the index-negative
subgroup.

Application to the pancreatic cancer trial was less conclusive and
illustrates the importance of separating numerical convergence from
scientific evidence. The three-biomarker model had satisfactory split
$\widehat R$ values, but all biomarker and interaction credible intervals
contained zero and the posterior treatment-effect intervals were wide. The
CA19-9/Axl analysis produced a subgroup with a more favorable descriptive
treatment comparison, consistent with the direction of earlier work
\cite{HeYe2018AstC}; however, its interaction credible interval also included
zero, and the index-negative subgroup contained only 74 patients. The other
pairwise scores generated either a very small subgroup or no division at all.
Accordingly, these data support the fitted scores as hypotheses for further
study rather than as validated treatment rules.

We note several limitations. First, the simulation studies used two
independent normally distributed biomarkers, exponential event times, modest
censoring, and treatment allocation of 0.5. They do not yet assess correlated
biomarkers, more than two biomarkers, alternative baseline hazards,
covariate-dependent censoring, or model misspecification. In addition,
simulation data sets were required to have at least 60 subjects in each
treatment-by-subgroup cell, and failed fits were retried until a successful
fit was obtained. The reported operating characteristics therefore apply to
eligible, successfully fitted data sets and should not be interpreted as
unconditional performance in sparse-subgroup settings.

Second, the subgroup indicator is a discontinuous function of
$\boldsymbol\gamma$. Small changes in an index coefficient can reclassify
subjects near the threshold, producing autocorrelation or multimodality even
when split $\widehat R$ is near one. The moderate effective sample sizes in
the application and the very small subgroup from some posterior-mean scores
are manifestations of this difficulty. Multiple dispersed starting values,
longer runs, sensitivity analyses for the prior variance, and posterior
summaries of each patient's subgroup-membership probability should accompany
future applications.

Third, the present analyses used $\lambda=0$ and therefore did not perform
formal biomarker selection. Although an $\ell_1$ penalty is available, it
shrinks coefficients without placing exact posterior mass at zero. Bayesian
spike-and-slab priors, continuous shrinkage priors, or a model-averaging
strategy may provide a more principled approach when many candidate
biomarkers are considered. Such extensions should be evaluated for both
selection accuracy and stability of the induced subgroup.

Finally, the Kaplan--Meier, log-rank, and conventional Cox interaction
analyses used subgroup assignments learned from the same trial. Their
p-values do not propagate uncertainty in $\boldsymbol\gamma$ and may be
optimistic. A confirmatory analysis should prespecify the biomarker score and
fixed threshold and then evaluate them in independent data. Internal
validation by sample splitting or nested resampling would provide a useful
intermediate assessment, but external validation remains necessary before a
score is used to guide treatment. Future work should also examine posterior
predictive measures of subgroup benefit that integrate over uncertain
membership rather than relying only on a hard classification from the
posterior-mean score.

Despite these limitations, the fixed-threshold parameterization provides a
simple and computationally practical framework for combining multiple
biomarkers. The simulation results establish encouraging estimation
properties, while the application demonstrates both the potential of the
CA19-9/Axl score and the uncertainty that must be addressed before the method
can support confirmatory biomarker-guided treatment decisions.
